\section{Operational Control and Resource Allocation}
\subsection{Decision Layers}
Once the priority field is known, the next question is not simply how much effort to spend, but how to structure that effort. We distinguish three control layers:
\begin{itemize}
  \item \textbf{anchoring}: repeatable patrol loops that guarantee physical presence in decisive cells;
  \item \textbf{persistence}: fixed devices that maintain continuous watch in stable hotspots;
  \item \textbf{adaptation}: drone sorties that absorb uncertainty and provide surge observation.
\end{itemize}
Those three layers become the decision variables of the controller:
\begin{itemize}
  \item \(x_r \in \mathbb{Z}_{+}\): number of times patrol loop \(r\) is activated;
  \item \(y_i \in \{0,1,2\}\): number of fixed devices assigned to zone \(i\);
  \item \(z_j \in \mathbb{Z}_{+}\): drone sorties assigned to aerial sector \(j\).
\end{itemize}

\subsection{Route Library and Patrol Geometry}
\subsubsection{Feasible Loop Construction}
Let \(\mathcal{G}_R=(V_R,E_R,w)\) again denote the weighted road graph. A patrol column \(r\) based at station \(s\in\mathcal{S}\) is encoded by edge-use variables \(\xi_e^{(r)}\in\{0,1\}\) and zone-coverage variables \(c_i^{(r)}\in\{0,1\}\). The route-generation subproblem is a resource-constrained prize-collecting tour:
\begin{align}
\max_{\xi^{(r)},c^{(r)}} \quad &
\Pi_r=\sum_{i=1}^{64} R_i c_i^{(r)}-\mu\sum_{e\in E_R} w_e \xi_e^{(r)}
\label{eq:routeobj}\\
\text{s.t.}\quad
& \sum_{e\in\delta^+(v)}\xi_e^{(r)}-\sum_{e\in\delta^-(v)}\xi_e^{(r)}=0,
\quad \forall v\in V_R\setminus\{s\},
\label{eq:routeflow}\\
& \sum_{e\in\delta^+(s)}\xi_e^{(r)}=
\sum_{e\in\delta^-(s)}\xi_e^{(r)}=1,
\label{eq:routeanchor}\\
& h_r=\sum_{e\in E_R} w_e \xi_e^{(r)}+t_{\text{stop}}\sum_{i=1}^{64} c_i^{(r)}\le h_{\max},
\label{eq:route}\\
& c_i^{(r)}\le \sum_{e\in\delta(i)}\xi_e^{(r)},\quad \forall i.
\label{eq:routecover}
\end{align}
Here \(\delta^+(v)\) and \(\delta^-(v)\) denote the outgoing and incoming road edges of node \(v\), \(\delta(i)\) denotes the set of road edges that enter zone \(i\), and standard subtour-elimination constraints are omitted for brevity. Equation~\eqref{eq:routeobj} makes the route library a genuine graph-theoretic object: each column is the solution of a bounded tour problem, not an arbitrary loop sketched after the fact.

The zone-coverage matrix is defined by
\begin{equation}
a_{ir} =
\begin{cases}
1, & \text{if } c_i^{(r)}=1,\\
0, & \text{otherwise}.
\end{cases}
\end{equation}
Patrol intensity in zone \(i\) is therefore
\begin{equation}
F_i = \sum_{r\in\mathcal{R}} a_{ir}x_r.
\end{equation}

\subsubsection{Enumeration and Pruning}
Equations~\eqref{eq:routeobj}--\eqref{eq:routecover} are the route-level version of team orienteering. In implementation, the master MILP does not solve that subproblem from scratch on every branch-and-bound node. Instead, for each station we generate a finite set of high-quality columns from \(K\)-shortest-path seeds and retain only the non-dominated loops. Route \(r\) dominates route \(r'\) if \(h_r\le h_{r'}\) and \(a_{ir}\ge a_{ir'}\) for every zone \(i\), with strict improvement in at least one component.

We additionally cap within-station overlap by
\begin{equation}
\operatorname{ov}(r,r')=
\frac{\sum_{e\in E_R}\xi_e^{(r)}\xi_e^{(r')}}{\min\left\{\sum_{e\in E_R}\xi_e^{(r)},\sum_{e\in E_R}\xi_e^{(r')}\right\}}
\le \rho_{\max},
\label{eq:overlap}
\end{equation}
and prune the resulting library by:
\begin{enumerate}
  \item removing any route that violates the time bound \eqref{eq:route};
  \item removing dominated routes whose coverage vector is inferior to a lower-cost alternative;
  \item limiting excessive overlap among routes from the same station so that the library preserves tactical variety.
\end{enumerate}
This produces a controller-ready route set that behaves like an operational menu rather than an arbitrary mathematical index. The route-cover matrix used in the code is therefore a compact column-generation interface to the underlying path-flow graph model.

\subsection{Protection Gain Model}
The controller values assets for their contribution to service intensity. Let
\[
\eta_i=\alpha F_i+\beta y_i+\gamma\sum_j b_{ij}z_j
\]
denote total surveillance intensity in zone \(i\), where \(b_{ij}=1\) if drone sector \(j\) covers zone \(i\). Zone-level protection gain is modeled as
\begin{equation}
E_i = 1-\exp\!\left(-\alpha F_i-\beta y_i-\gamma \sum_j b_{ij}z_j\right).
\label{eq:gain}
\end{equation}
This is a deliberately saturating response: the first patrol, first device, and first drone window matter much more than later increments.

The gain model can be derived from an event-level view. If illegal arrivals in zone \(i\) are Poisson with mean \(\lambda_i\), then expected prevented weighted loss in one planning period is
\begin{equation}
\mathbb{E}[\Delta_i]=R_i\lambda_i\left(1-e^{-\eta_i}\right)\bar{g}_i,
\label{eq:delta}
\end{equation}
where \(\bar{g}_i\) is expected interception success conditional on detection.

\begin{proposition}
For fixed \(\bar{g}_i\), the protection-gain function \(E_i(\eta_i)=1-e^{-\eta_i}\) is increasing and strictly concave.
\end{proposition}
\begin{proof}
\[
\frac{\partial E_i}{\partial \eta_i}=e^{-\eta_i}>0,\qquad
\frac{\partial^2 E_i}{\partial \eta_i^2}=-e^{-\eta_i}<0.
\]
The controller therefore prefers concentrated reinforcement in decisive cells but recognizes diminishing marginal returns.
\end{proof}

\subsection{MILP Approximation and Control Insight}
Because the solver is mixed-integer linear, we use a capped linear surrogate:
\[
e_i = \min\{1,\alpha F_i+\beta y_i+\gamma\sum_j b_{ij}z_j\}.
\]
The implementation enforces
\begin{align}
e_i &\le \alpha F_i+\beta y_i+\gamma\sum_j b_{ij}z_j, \label{eq:lin1}\\
e_i &\le 1. \label{eq:lin2}
\end{align}
which preserves monotonicity and tractability while remaining close to the original concave response.

Ignoring integrality and hotspot floors, the continuous relaxation equalizes weighted marginal value across active zones:
\begin{equation}
R_i e^{-\eta_i} \approx R_k e^{-\eta_k}.
\label{eq:balancing}
\end{equation}
Equation~\eqref{eq:balancing} is the core allocation insight. Scarce resources are pulled toward cells with large \(R_i\) until their marginal value is driven down to the same level as other active cells. That is why the optimized strategy concentrates along the eastern--southern corridor instead of dispersing effort uniformly.

\subsection{Controller Objective and Constraints}
The patrol component is a path-flow master problem written in route columns. Each integer variable \(x_r\) is the multiplicity of one graph-feasible patrol tour, so connectivity and return-to-station feasibility are already enforced upstream by Equations~\eqref{eq:routeflow}--\eqref{eq:routecover}. The master problem only decides how often each feasible column should be activated alongside devices and UAVs.

The optimization model maximizes weighted expected protection net of route burden:
\begin{equation}
\max \; Z = \sum_{i=1}^{64} R_iE_i - \lambda \sum_{r\in\mathcal{R}} h_rx_r,
\label{eq:obj}
\end{equation}
with \(E_i\) replaced by the linearized surrogate \(e_i\) in implementation.

The controller is subject to five feasibility rules:
\begin{align}
\sum_{r\in\mathcal{R}} 2h_r x_r &\le H_R, \label{eq:rangerhours}\\
\sum_{i=1}^{64} y_i &\le M, \label{eq:devices}\\
\sum_j z_j &\le U, \label{eq:drones}\\
\sum_{r: s(r)=s} x_r &\le \kappa_s, \quad \forall s, \label{eq:station}\\
F_i + y_i + \sum_j b_{ij}z_j &\ge \underline{q}\,\mathbb{I}(R_i\ge q_{0.75}),
\label{eq:hotspots}
\end{align}
where \(H_R\) is total ranger-hours, \(M\) is device inventory, \(U\) is available drone sorties, and \(\kappa_s\) is station capacity. The hotspot-floor constraint \eqref{eq:hotspots} is the mathematical expression of the operating doctrine: decisive cells must not be allowed to fall below a minimum service level.

\subsection{Staffing Floor and Rolling-Horizon Control}
If \(\mathcal{H}\) is the hotspot set and \(f_i^\star\) is the required revisit frequency in those cells, then a conservative staffing floor is
\begin{equation}
n_{\min} \ge \frac{\sum_{i\in\mathcal{H}} f_i^\star \bar{h}_i}{h_{\text{shift}}\rho},
\label{eq:rangermin2}
\end{equation}
where \(\bar{h}_i\) is the representative route-hours required to protect hotspot \(i\), \(h_{\text{shift}}\) is usable shift time, and \(\rho\) is field utilization. This equation closes the loop between hotspot protection and staffing demand.

The controller runs in a rolling horizon:
\begin{table}[H]
\centering
\caption{Closed-loop protection controller.}
\label{alg:planner}
\begin{tabular}{p{0.95\linewidth}}
\toprule
\textbf{Input:} zone features \((W_i,S_i,H_i,A_i,D_i)\), route library \(\mathcal{R}\), capacities \((H_R,M,U)\), and horizon \(T\). \\
\textbf{Output:} deployment plan \(\{x_r^t,y_i^t,z_j^t\}_{t=1}^T\) and staffing floor \(n_{\min}\). \\
\midrule
1. Build the priority field \(R_i\) from Equations~\eqref{eq:risk}--\eqref{eq:prior}. \\
2. Generate graph-feasible patrol columns from Equations~\eqref{eq:routeobj}--\eqref{eq:routecover} and build \(a_{ir}\) and \(b_{ij}\). \\
3. Solve the allocation controller under Equations~\eqref{eq:obj}--\eqref{eq:hotspots}. \\
4. Stress test the resulting plan through detection, dispatch, and interception simulation. \\
5. Update patrol pressure \(C_i\), refresh \(P_i\), and re-optimize if performance falls below threshold. \\
6. Return the controller output and the minimum staffing recommendation. \\
\bottomrule
\end{tabular}
\end{table}
Figure~\ref{fig:controllerarch} expands the same logic visually by separating the input stack, the optimization core, and the deployable products that leave the solver.

\begin{figure}[H]
  \centering
  \includegraphics[width=\linewidth]{fig16_controller_architecture.png}
  \caption{Algorithm architecture of the allocation controller, showing how state estimates, feasible-route geometry, and resource budgets pass through the MILP core to produce patrol, sensor, and UAV tasking.}
  \label{fig:controllerarch}
\end{figure}

\subsection{Baseline Control Package}
Under baseline assumptions, the controller assigns 49\% of patrol-hours to the eastern sector, 31\% to the southern--central corridor, and 20\% to western and northern sectors. Eight of the twelve fixed devices are anchored in the top-risk quartile, while four drone sorties are rotated across three high-uncertainty sectors. This is not a cosmetic rebalance; it is a structural shift from thin park-wide patrol coverage to corridor-based service assurance.

\begin{figure}[H]
  \centering
  \includegraphics[width=\linewidth]{fig05_patrol_routes.png}
  \caption{Controller-selected patrol families over the official Etosha basemap, showing how route geometry is concentrated into the decisive corridor.}
  \label{fig:routes}
\end{figure}

\begin{figure}[H]
  \centering
  \includegraphics[width=\linewidth]{fig06_allocation_dashboard.png}
  \caption{Control ledger translating manpower, devices, and drone effort into sector-level service commitments.}
  \label{fig:allocation}
\end{figure}

\begin{figure}[H]
  \centering
  \includegraphics[width=\linewidth]{fig13_sensor_drone_map.png}
  \caption{Layered control design on a real-map GIS basemap, in which fixed sensors hold persistent hotspots and UAVs absorb uncertainty and surge demand.}
  \label{fig:techmap}
\end{figure}
\FloatBarrier
